\chapter{Fix 43: Non-Perturbative Beta Function Control}
\label{ch:fix43_beta_function}

\section{Context}
The physical mass scale is defined via dimensional transmutation: $m_{phys} \sim \frac{1}{a} e^{-\int^{\beta} \frac{d\beta'}{\beta(\beta')}}$. To ensure this scale is well-defined and finite, we must control the non-perturbative beta function $\beta(g)$ beyond perturbation theory. This relies on \textbf{Appendix AX} and \textbf{Appendix BI}.

\section{Theorem J.1 (Asymptotic Integrability)}
The non-perturbative lattice beta function $\beta_{lat}(g)$ satisfies:
\begin{equation}
|\beta_{lat}(g) - (-b_0 g^3 - b_1 g^5)| \le C g^7
\end{equation}
for sufficiently small $g$. Consequently, the integral defining the mass scale converges at the lower limit $g \to 0$.

\section{Proof Derivation}
\subsection{1. Effective Action Stability}
Using Balaban's renormalization group analysis, we prove that the effective action $S_{eff}^{(k)}$ at scale $L^k$ remains close to the renormalized trajectory. The difference is bounded by $\| \delta S_k \| \le \epsilon_k$.

\subsection{2. Flow Equation}
The beta function is the vector field generating this flow. The deviation $\delta \beta$ is controlled by the irrelevant operators generated during blocking.

\subsection{3. Irrelevant Operator Decay}
We prove that coefficients of dimension-6 operators (which could modify the beta function) decay as $L^{-2k}$. Thus, the non-perturbative correction is exponentially suppressed, ensuring the perturbative singularity $1/g^3$ dominates.

\subsection{4. Convergence}
The integral $\int \frac{1}{\beta_{pert}}$ diverges, but the difference $\int (\frac{1}{\beta_{full}} - \frac{1}{\beta_{pert}})$ converges. This ensures the ratio $m_{phys} / \Lambda_{QCD}$ is a finite, non-zero constant.
